\section{Introduction}

The United States interstate highway system was designed for an era of stable climate. The
engineering specifications that governed Eisenhower-era construction --- drainage sizing, bridge
freeboard, roadway elevation above ordinary high water --- reflected the precipitation and sea
level conditions of the mid-twentieth century. In the intervening seventy years, both the climate
envelope and the infrastructure's exposure within it have shifted materially. FHWA's current
investment prioritization frameworks --- the National Highway Performance Program, the PROTECT
climate resilience grant program, and the Surface Transportation Block Grant --- do not include
a standardized climate exposure criterion for highway corridors. Funding decisions that should
incorporate climate-driven closure risk instead use proxy metrics (total SFHA miles, V/C ratio)
that systematically undervalue structural vulnerability.

This paper addresses that gap. We introduce the D1 dimension of the ROUTE rubric: a composite
climate vulnerability score, 0--10, computed for all 227 interstate corridors in the ROUTE
v1.2 corpus. D1 captures two distinct exposure mechanisms. For flood-prone corridors, D1
measures the length of highway within FEMA Special Flood Hazard Areas (SFHA), with additional
weight for \emph{consecutive} SFHA miles --- the condition in which a road remains inside the
100-year flood zone without interruption, eliminating internal escape routes within the flooded
reach. For winter precipitation and wildfire corridors, D1 translates closure frequency and
duration into an equivalent vulnerability score through an empirical calibration anchored on
Donner Pass (I-80) as the reference worst-case winter corridor.

Three findings anchor this paper. First, I-10 through Gulf Coast Louisiana is the most
flood-exposed corridor in the national system: 127 consecutive SFHA miles from Lake Charles
to New Orleans, a mean elevation of 0.8 meters above sea level, and a subsidence rate of
1--2 cm per year that compounds exposure beyond what static flood maps capture. No other
interstate corridor combines this length of consecutive inundation exposure with this
proximity to sea level. Second, I-80 Donner Pass is the highest winter-closure corridor in
the corpus, scoring D1 = 7.8 on closure frequency and duration despite no flood exposure ---
demonstrating that the D1 dimension is not synonymous with flood risk. Third, under NOAA
intermediate sea level rise projections (0.5m by 2050), the Gulf Coast I-10 D1 score
increases from 8.4 to a projected 9.1, shifting the national investment priority order:
the 2050 profile places Gulf Coast I-10 hardening as the single most urgent climate
adaptation investment in the highway system, ahead of Donner snowshed construction and ahead
of the current PROTECT program's funding targets.

The policy implication is direct. FHWA's PROTECT program (\$8.7 billion over five years
under the Infrastructure Investment and Jobs Act) currently prioritizes corridors by total
SFHA miles. The consecutive-miles metric introduced here identifies Gulf Coast Louisiana
I-10 as uniquely exposed in a way that the total-miles criterion does not capture --- because
the consecutive nature of the SFHA exposure means that when the corridor floods, there is no
segment of the same road available as an internal alternate. Correcting this metric would
shift PROTECT funding toward the highest-consequence corridors.

\subsection*{Contributions}

\begin{enumerate}
\item A climate exposure score (D1) for all 227 interstate corridors, combining FEMA SFHA
flood overlay and historical closure frequency data under a single calibrated rubric
\item A demonstration that consecutive SFHA miles, rather than total SFHA miles, is the
structurally correct exposure metric --- and that this distinction changes the investment
priority ranking
\item Quantified 2050 D1 projections for the top-10 exposure corridors under NOAA intermediate
and high sea level rise scenarios, producing a mid-century investment priority order
\item A policy recommendation: FHWA PROTECT program criteria should incorporate consecutive
SFHA miles and 2050 projected exposure, not only current total SFHA mileage
\end{enumerate}
